%!TEX root = ../../main.tex
%% ===============================================================
%% 2.1 리그 오브 레전드와 Match-V5 텔레메트리
%% 파일: chapters/02_background/01_telemetry.tex
%% 상위: chapters/02_background/chapter.tex
%% 정의 라벨: sec:bg-telemetry
%% 참조 라벨: -
%% 포함 객체: -
%% 인용: -
%% 상태: 초안 (docs/OUTLINE.md 참고)
%% ===============================================================

\section{리그 오브 레전드와 Match-V5 텔레메트리}
\label{sec:bg-telemetry}

LoL의 한 경기는 블루 팀(team 100)과 레드 팀(team 200) 각 5 명, 총 $N=10$ 명의 참가자로 구성된다. 각 참가자는 경기 시작 전 챔피언 하나, 소환사 주문 둘, 룬 세트를 선택하며 경기 중 골드로 아이템을 구매한다. 지도 좌표계는 좌하단을 원점으로 하는 $[0, 16{,}000]^2$ 정수 격자이다.

Riot Match-V5 API는 경기당 두 개의 JSON 문서를 제공한다. \textbf{detail}은 경기 메타데이터(패치 버전, 참가자·챔피언·룬·최종 통계)를, \textbf{timeline}은 시계열을 담는다. timeline은 (i) 60{,}000 ms 간격의 \emph{프레임}(frame) 열과 (ii) 밀리초 타임스탬프를 갖는 \emph{사건}(event) 열로 구성된다. 프레임에는 참가자별 위치 $(x, y)$, 레벨, 경험치, 현재·누적 골드, 미니언 처치 수, 체력·자원, 챔피언 스탯(공격력, 방어력 등 25 종), 피해 통계(12 종)가 들어 있다. 사건에는 \texttt{CHAMPION\_KILL}, \texttt{CHAMPION\_SPECIAL\_KILL} (에이스, 다중 킬, 첫 킬), \texttt{ELITE\_MONSTER\_KILL}, \texttt{BUILDING\_KILL}, \texttt{TURRET\_PLATE\_DESTROYED}, \texttt{WARD\_PLACED}, \texttt{WARD\_KILL}, \texttt{ITEM\_PURCHASED} 등이 있으며, 킬 사건은 위치 좌표와 가해자·피해자·어시스트 참가자 ID를 포함한다.

이 구조가 본 논문의 모든 설계를 규정한다. 위치와 상태는 분 단위로만 알 수 있으므로 교전의 시공간 경계는 사건(특히 킬)으로부터 추론해야 하고, 모델의 입력은 60 초 프레임을 어떤 방식으로든 재표본화한 것일 수밖에 없다.
